\section{The incidence graph and the row graph are not substitutes}
\label{sec:two-graphs}

The completion incidence graph of TPC-77/78 is bipartite:
\[
 G_C\subseteq I_{\rm analysis}\sqcup J_{\rm variables}.
\]
It factorizes or solves the coefficient program
\(\min_\theta\|b+C\theta\|_1\).  The present graph
\[
 G_{\rm row}(P)\subseteq V_{\rm physical}\times V_{\rm physical}
\]
controls the Perron eigenvector and scalar Hermitian phase transport.

\begin{proposition}[Graph non-substitution]\label{prop:non-substitution}
Neither of the following implications holds without an additional
literal construction:
\[
 G_C\text{ is a forest}
 \ \Longrightarrow\
 G_{\rm row}(P)\text{ is a forest},
\]
\[
 G_{\rm row}(P)\text{ is a forest}
 \ \Longrightarrow\
 G_C\text{ is a forest}.
\]
\end{proposition}

\begin{proof}
Take \(C=I_4\).  Its bipartite support is a matching, hence a forest,
while an independently assembled row majorant may be the adjacency
matrix of \(C_4\), whose row graph has a cycle.

Conversely, take a \(2\times2\) matrix \(C\) with all four entries
nonzero.  Its support is \(K_{2,2}\), which has a cycle, while a
two-row Hermitian majorant with one nonzero off-diagonal entry has a
tree row graph.  These examples can share the same analysis-row labels;
support of one matrix does not determine support of the other.
\end{proof}

\begin{remark}[Different phase transports]
TPC-78 transports a dual charge through degree-two variable nodes of a
bipartite tree.  TPC-80 switches eigenvector phases along Hermitian row
edges.  Their equations, vertices, and obstruction spaces are
different even when both happen to be one-dimensional.
\end{remark}

The literal interface must therefore archive both graphs, with distinct
names and hashes.  A cycle counted before post-\(q\) entry aggregation
is not a row holonomy cycle if its final Hermitian edge cancels to zero.
